% Shared preamble of the proof overview documents  <Folder>/<Name>_overview.tex.
% Each overview starts with
%     % Shared preamble of the proof overview documents  <Folder>/<Name>_overview.tex.
% Each overview starts with
%     % Shared preamble of the proof overview documents  <Folder>/<Name>_overview.tex.
% Each overview starts with
%     \input{../../scripts/overview_preamble.tex}
% (for ThinMatching/: \input{../scripts/overview_preamble.tex}) and is compiled from its own folder by
% scripts/build_overviews.py.
\documentclass[10pt,a4paper]{article}
\usepackage[T1]{fontenc}
\usepackage[utf8]{inputenc}
\usepackage[margin=2.2cm]{geometry}
\usepackage{amsmath,amssymb,amsthm}
\usepackage{array,booktabs,longtable}
\usepackage[dvipsnames]{xcolor}
\usepackage{enumitem}
\usepackage{tikz}
\usetikzlibrary{arrows.meta,positioning,shapes.multipart,fit,calc}
\usepackage[hidelinks]{hyperref}
\setlist{itemsep=2pt,topsep=4pt}
\renewcommand{\arraystretch}{1.2}

% Lean names, files and paper references
\newcommand{\lean}[1]{\textcolor{blue!50!black}{\texttt{\detokenize{#1}}}}
\newcommand{\file}[1]{\textcolor{OliveGreen!70!black}{\texttt{\detokenize{#1}}}}
\newcommand{\paperref}[1]{\textsf{#1}}
\newcommand{\hyp}[1]{\textcolor{BrickRed}{\textbf{#1}}}

% Boxes for the dependency diagrams
\tikzset{
  res/.style={draw, rounded corners, fill=blue!6, align=center, inner sep=4pt, font=\small},
  main/.style={res, fill=blue!18, very thick},
  cited/.style={res, fill=red!10, dashed},
  shared/.style={res, fill=gray!12},
  dep/.style={-{Stealth[length=2.2mm]}, thick},
}

\newcommand{\overviewtitle}[2]{%
  \begin{center}{\Large\bfseries #1}\\[4pt]{\small Proof overview of the Lean~4/Mathlib formalisation}\\[2pt]
  {\small #2}\end{center}}

% Standard paragraph describing the paper
\newcommand{\paperversion}{\noindent\textbf{Paper.} J.~Breitung, T.~Gervens, M.~Grohe.
  \emph{Graph Similarity Meets Optimal Transport}.\\ File \file{Graph_Similarity_Meets_Optimal_Transport.pdf} in the
  repository; statement, equation, page and line numbers refer to this PDF. \file{docs/NUMBERING.md} lists the Lean
  files and declarations of all numbered statements.\par\smallskip}

% (for ThinMatching/: \input{../scripts/overview_preamble.tex}) and is compiled from its own folder by
% scripts/build_overviews.py.
\documentclass[10pt,a4paper]{article}
\usepackage[T1]{fontenc}
\usepackage[utf8]{inputenc}
\usepackage[margin=2.2cm]{geometry}
\usepackage{amsmath,amssymb,amsthm}
\usepackage{array,booktabs,longtable}
\usepackage[dvipsnames]{xcolor}
\usepackage{enumitem}
\usepackage{tikz}
\usetikzlibrary{arrows.meta,positioning,shapes.multipart,fit,calc}
\usepackage[hidelinks]{hyperref}
\setlist{itemsep=2pt,topsep=4pt}
\renewcommand{\arraystretch}{1.2}

% Lean names, files and paper references
\newcommand{\lean}[1]{\textcolor{blue!50!black}{\texttt{\detokenize{#1}}}}
\newcommand{\file}[1]{\textcolor{OliveGreen!70!black}{\texttt{\detokenize{#1}}}}
\newcommand{\paperref}[1]{\textsf{#1}}
\newcommand{\hyp}[1]{\textcolor{BrickRed}{\textbf{#1}}}

% Boxes for the dependency diagrams
\tikzset{
  res/.style={draw, rounded corners, fill=blue!6, align=center, inner sep=4pt, font=\small},
  main/.style={res, fill=blue!18, very thick},
  cited/.style={res, fill=red!10, dashed},
  shared/.style={res, fill=gray!12},
  dep/.style={-{Stealth[length=2.2mm]}, thick},
}

\newcommand{\overviewtitle}[2]{%
  \begin{center}{\Large\bfseries #1}\\[4pt]{\small Proof overview of the Lean~4/Mathlib formalisation}\\[2pt]
  {\small #2}\end{center}}

% Standard paragraph describing the paper
\newcommand{\paperversion}{\noindent\textbf{Paper.} J.~Breitung, T.~Gervens, M.~Grohe.
  \emph{Graph Similarity Meets Optimal Transport}.\\ File \file{Graph_Similarity_Meets_Optimal_Transport.pdf} in the
  repository; statement, equation, page and line numbers refer to this PDF. \file{docs/NUMBERING.md} lists the Lean
  files and declarations of all numbered statements.\par\smallskip}

% (for ThinMatching/: \input{../scripts/overview_preamble.tex}) and is compiled from its own folder by
% scripts/build_overviews.py.
\documentclass[10pt,a4paper]{article}
\usepackage[T1]{fontenc}
\usepackage[utf8]{inputenc}
\usepackage[margin=2.2cm]{geometry}
\usepackage{amsmath,amssymb,amsthm}
\usepackage{array,booktabs,longtable}
\usepackage[dvipsnames]{xcolor}
\usepackage{enumitem}
\usepackage{tikz}
\usetikzlibrary{arrows.meta,positioning,shapes.multipart,fit,calc}
\usepackage[hidelinks]{hyperref}
\setlist{itemsep=2pt,topsep=4pt}
\renewcommand{\arraystretch}{1.2}

% Lean names, files and paper references
\newcommand{\lean}[1]{\textcolor{blue!50!black}{\texttt{\detokenize{#1}}}}
\newcommand{\file}[1]{\textcolor{OliveGreen!70!black}{\texttt{\detokenize{#1}}}}
\newcommand{\paperref}[1]{\textsf{#1}}
\newcommand{\hyp}[1]{\textcolor{BrickRed}{\textbf{#1}}}

% Boxes for the dependency diagrams
\tikzset{
  res/.style={draw, rounded corners, fill=blue!6, align=center, inner sep=4pt, font=\small},
  main/.style={res, fill=blue!18, very thick},
  cited/.style={res, fill=red!10, dashed},
  shared/.style={res, fill=gray!12},
  dep/.style={-{Stealth[length=2.2mm]}, thick},
}

\newcommand{\overviewtitle}[2]{%
  \begin{center}{\Large\bfseries #1}\\[4pt]{\small Proof overview of the Lean~4/Mathlib formalisation}\\[2pt]
  {\small #2}\end{center}}

% Standard paragraph describing the paper
\newcommand{\paperversion}{\noindent\textbf{Paper.} J.~Breitung, T.~Gervens, M.~Grohe.
  \emph{Graph Similarity Meets Optimal Transport}.\\ File \file{Graph_Similarity_Meets_Optimal_Transport.pdf} in the
  repository; statement, equation, page and line numbers refer to this PDF. \file{docs/NUMBERING.md} lists the Lean
  files and declarations of all numbered statements.\par\smallskip}

\begin{document}
\overviewtitle{Theorem 15: NP-hardness of the optimal-transport edit distance}{Folder \file{GsimOT/Theorem15_EditHardness}}
\paperversion

\section{Statement}

\paragraph{Paper.} Theorem~15 (Section~4, p.~12, l.~418; proof in Appendix~G, pp.~50--51, l.~1577--1620):
\emph{$\textsc{GSim}(d^{\mathrm{OT}}_{\mathrm{ed}})$ is NP-hard, even when restricted to collections of paths.}
It is part~(1) of Theorem~4 (p.~4, l.~134). The proof is the reduction of [GRW18, Theorem~8] from the strongly
NP-hard problem Three-Partition: an instance $(A,a_1,\dots,a_{3m})$ in unary with $\sum_i a_i=mA$ and
$A/4<a_i<A/2$ is mapped to $(F_1,F_2,t)$, where $F_1$ is the disjoint union of paths on $a_i$ vertices,
$F_2$ the disjoint union of $m$ paths on $A$ vertices, and $t=4m$. The claim (p.~50, l.~1584--1585) is that
$d^{\mathrm{OT}}_{\mathrm{ed}}(F_1,F_2)\le t$ iff the instance is a YES-instance.

\paragraph{Lean.} Only the correctness of the reduction is formalised. With
$F_1=$ \lean{pathUnion I.a} and $F_2=$ \lean{pathUnion fun _ : Fin I.m => I.A}:
\begin{itemize}
\item \lean{GsimOT.Hardness.ThreePartition.dOTEdit_le_iff}: if $a_i\ge2$ for all $i$, then
  \[ d^{\mathrm{OT}}_{\mathrm{ed}}(F_1,F_2)\le 4m \iff I \text{ is a YES-instance.} \]
\item \lean{GsimOT.Hardness.ThreePartition.dOTEdit_shift_le_iff}: for \emph{every} instance $I$, the same
  equivalence for the shifted instance (add $3$ to each $a_i$ and $9$ to $A$, l.~1586--1587), whose
  YES/NO status equals that of $I$.
\end{itemize}
Here $d^{\mathrm{OT}}_{\mathrm{ed}}(G,H)=\min_{Q\in\mathrm{DS}(V,W)}\|B(G,H,Q)\|_{(1)}$ is \lean{GsimOT.dOTEdit}
(\file{Common/EditDistance.lean}), the same definition used by Proposition~5 and Theorem~6; a
Three-Partition instance is the structure \lean{GsimOT.Hardness.ThreePartition} (fields \lean{sum_eq},
\lean{lt_four_mul}: $A<4a_i$, \lean{two_mul_lt}: $2a_i<A$), and \lean{IsYes} asks for
$\sigma\colon[3m]\to[m]$ with classes of size $3$ and sum $A$.

\section{Where the proof lives}
\begin{longtable}{>{\raggedright\arraybackslash}p{5.2cm} >{\raggedright\arraybackslash}p{10.6cm}}
\toprule \textbf{File} & \textbf{Content} \\ \midrule\endhead \bottomrule\endfoot
\file{Theorem15_EditHardness/Theorem15.lean} & $\mathrm{lost}(Q)$, the identity $\|B\|_{(1)}=4m+2\,\mathrm{lost}(Q)$,
  both directions of the reduction, the shifted version. \\
\file{Common/ThreePartition.lean} & Three-Partition instances, the layout bijection of a solution
  (\lean{layoutEquiv}), the counting step \lean{isYes_of_isDS}, the shift. \\
\file{Common/PathsAndCycles.lean} & \lean{pathUnion}, edge count \lean{sum_adj_pathUnion}, neighbours
  (\lean{exists_pathUnion_adj}), constancy along paths (\lean{eq_of_pathUnion_adj}). \\
\file{Common/EditDistance.lean}, \file{Common/MatrixNorms.lean} & \lean{dOTEdit}, \lean{dOTEdit_le},
  attainment \lean{exists_dOTEdit_eq}, $C_{\mathrm{ed}}(Q)=\|B(G,H,Q)\|_{(1)}$ (\lean{costEdit_eq_entrywiseNorm_B}). \\
\end{longtable}

\section{Proof layout}
\begin{center}
\begin{tikzpicture}
  \node[main] (thm) at (0,0) {Theorem 15 (any instance)\\\lean{dOTEdit_shift_le_iff}};
  \node[shared] (shift) at (6,0) {shift\\\lean{shift_isYes_iff}};
  \node[main] (red) at (0,-1.8) {Theorem 15, $a_i\ge2$\\\lean{dOTEdit_le_iff}};
  \node[res] (cost) at (-5.4,-4) {$\|B\|_{(1)}=4m+2\,\mathrm{lost}(Q)$\\\lean{entrywiseNorm_B_eq}\\\lean{sum_adj_sub_sum_adj}};
  \node[res] (zero) at (0,-4) {$\mathrm{lost}(Q)=0$ $\Rightarrow$ no edge\\to a non-edge\\\lean{mul_eq_zero_of_lost_eq_zero}};
  \node[res] (perm) at (5.4,-4) {YES: the layout has\\$\mathrm{lost}=0$\\\lean{lost_permMat}};
  \node[shared] (att) at (-5.4,-6.3) {minimum attained\\\lean{exists_dOTEdit_eq}};
  \node[shared] (count) at (0,-6.3) {counting step\\\lean{isYes_of_isDS}};
  \node[shared] (lay) at (5.4,-6.3) {layout bijection\\\lean{layoutEquiv}\\\lean{layout_adj_pathUnion}};
  \draw[dep] (thm) -- (red); \draw[dep] (thm) -- (shift);
  \draw[dep] (red) -- (cost.north); \draw[dep] (red) -- (zero); \draw[dep] (red) -- (perm.north);
  \draw[dep, rounded corners] (red.west) -- (-7.9,-1.8) -- (-7.9,-6.3) -- (att.west);
  \draw[dep] (zero) -- (count); \draw[dep] (perm) -- (lay);
\end{tikzpicture}
\end{center}
Blue: this folder; grey: \file{GsimOT/Common}. Nothing is cited.

\section{The proof step by step}

\paragraph{Step 0: w.l.o.g.\ $a_i\ge4$ (l.~1586--1587).} The paper adds $3$ to every $a_i$ and $9$ to $A$.
Lean: \lean{ThreePartition.shift}, \lean{four_le_shift_a}, and \lean{shift_isYes_iff} (each class of a
solution has exactly three elements, \lean{card_eq_three_of_sum_eq}, so the class sums shift by $9$).
The formalised reduction \lean{dOTEdit_le_iff} assumes $a_i\ge2$, so that every vertex of $F_1$ has a
neighbour (\lean{exists_pathUnion_adj}).

\paragraph{Step 1: passing to blowups (l.~1585--1586).} The paper uses
$d^{\mathrm{OT}}_{\mathrm{ed}}(F_1,F_2)=d^{\bullet}_{\mathrm{ed}}(F_1,F_2)$, cited from [Gro25, Theorem~3.10],
and argues with alignments of blowups. \emph{The formalisation skips this step} and works with
$Q\in\mathrm{DS}(V,W)$ directly; see Section~6.

\paragraph{Step 2: YES-instances (l.~1588--1597).} A solution $\sigma$ lays the paths of $F_1$ out
consecutively on the paths of $F_2$; the bijection $\rho$ maps edges to edges, and only the $2m$ joining
edges of $F_2$ are not hit, so the cost is $4m$ (the paper then lifts $\rho$ to all blowups).
Lean: \lean{layoutEquiv} and \lean{layout_adj_pathUnion} (Common), \lean{lost_permMat}
($\mathrm{lost}(P_\rho)=0$), and Step~3 gives $C_{\mathrm{ed}}(P_\rho)=4m$, hence
$d^{\mathrm{OT}}_{\mathrm{ed}}\le4m$ by \lean{dOTEdit_le}.

\paragraph{Step 3: the edge count (p.~51, l.~1600--1602).} The paper counts the number $b$ of edges of $F_1'$ that
an alignment $\pi$ of the $k$-blowups maps to non-edges and obtains
$\|A^{\pi}_{F_1'}-A_{F_2'}\|_{(1)}=4mk^2+4b$. Lean, for every $Q\in\mathrm{DS}$: with
\[ \mathrm{lost}(Q)=\sum_{(v,w),(v',w')}Q(v,w)\,A_{F_1}(v,v')\,(1-A_{F_2}(w,w'))\,Q(v',w')\ \ge0 \]
(\lean{lost}, \lean{lost_nonneg}), the identity $|a-b|=b-a+2a(1-b)$ for $a,b\in\{0,1\}$ and the row and
column sums of $Q$ give
\[ \|B(F_1,F_2,Q)\|_{(1)}=\textstyle\sum A_{F_2}-\sum A_{F_1}+2\,\mathrm{lost}(Q)=4m+2\,\mathrm{lost}(Q) \]
(\lean{entrywiseNorm_B_eq}, \lean{sum_adj_sub_sum_adj}). This is the paper's identity with
$b/k^2$ replaced by $\mathrm{lost}(Q)/2$ ($\mathrm{lost}$ counts ordered pairs).

\paragraph{Step 4: NO-direction, clouds go into one path (l.~1598--1616).} The paper sets
$\gamma=1/(2mA)$, takes an alignment of cost at most $(4m+4\gamma)k^2$, deduces $b\le\gamma k^2$, so that
each cloud has a unique target path $j(v)$ receiving $\ge(1-\gamma)k$ of its vertices, and that $j(\cdot)$ is
constant along edges. Lean: since the minimum is attained (\lean{exists_dOTEdit_eq}), an optimal $Q$ has
cost $\le4m$, hence $\mathrm{lost}(Q)=0$ exactly; so $Q(v,w)Q(v',w')=0$ whenever $vv'\in E(F_1)$ and
$ww'\notin E(F_2)$ (\lean{mul_eq_zero_of_lost_eq_zero}). As every $v$ has a neighbour, the whole row
of $v$ lies in one path of $F_2$, the same for adjacent vertices, hence for all vertices of a path of $F_1$
(\lean{eq_of_pathUnion_adj}). These steps are inside the proof of \lean{dOTEdit_le_iff}
(\lean{hcomp}, \lean{hc}, \lean{hconst}).

\paragraph{Step 5: counting (l.~1617--1620).} The paper shows $|A-S_j|<\gamma mA=\tfrac12$, so $S_j=A$ and,
as $A/4<a_i<A/2$, exactly three paths go to path $j$. Lean: \lean{isYes_of_isDS} sums the column sums of
$Q$ over path $j$ of $F_2$ and gets $S_j=A$ exactly; \lean{card_eq_three_of_sum_eq} gives the triples.

\section{Assumptions and trusted base}
No hypotheses: the result is unconditional and uses only the axioms \lean{propext},
\lean{Classical.choice}, \lean{Quot.sound} (checked by \file{scripts/CheckAxioms.lean}). The complexity side
(polynomial size of the construction, strong NP-hardness of Three-Partition [GJ79]) is not formalised.
Proposition~57 (p.~57, l.~1805--1808: membership of $\textsc{GSim}(\delta^{\mathrm{OT}}_{\mathrm{ed}})$ in NP
and exact finite blowups) is not formalised.

\section{Notes on the formalisation}
\begin{itemize}
\item \textbf{No blowups, nothing cited.} The paper argues with blowups and needs
  $d^{\mathrm{OT}}_{\mathrm{ed}}=d^{\bullet}_{\mathrm{ed}}$ from [Gro25, Theorem~3.10] (p.~50, l.~1585--1586).
  Theorem~1 of the paper, which is formalised
  (\lean{GsimOT.theorem1}), covers only the operator norms, not the edit distance, so following the paper
  would introduce a cited hypothesis. The Lean proof is the fractional form of the same counting
  ($b/k^2\leftrightarrow\mathrm{lost}(Q)/2$, the cloud fractions $n_j(v)/k\leftrightarrow$ the row of $Q$),
  and the $\gamma$-approximation becomes exact because the minimum over $\mathrm{DS}$ is attained.
\item \textbf{Hypotheses.} \lean{dOTEdit_le_iff} assumes $a_i\ge2$ for all $i$;
  \lean{dOTEdit_shift_le_iff} applies it to the shifted instance (Step~0) and holds for every instance.
\item The Lean names \lean{dOTEdit} and \lean{costEdit} stand for the paper's $d^{\mathrm{OT}}_{\mathrm{ed}}$ and
  $C_{\mathrm{ed}}$.
\end{itemize}

\section{How to check this result}
From the repository root (after \texttt{export PATH=\$HOME/.elan/bin:\$PATH}):
\begin{itemize}
\item build the module:\\ \texttt{lake build GsimOT.Theorem15\_EditHardness.Theorem15}\\
  (or, after a full build, \texttt{lake env lean GsimOT/Theorem15\_EditHardness/Theorem15.lean});
\item check the axioms: \texttt{bash build.sh} runs \file{scripts/CheckAxioms.lean}, which contains\\
  \texttt{\#print axioms GsimOT.Hardness.ThreePartition.dOTEdit\_le\_iff}\\
  \texttt{\#print axioms GsimOT.Hardness.ThreePartition.dOTEdit\_shift\_le\_iff};
\item the table \file{Theorem15_correspondence.pdf} in this folder lists every paper item with its Lean
  declaration and source line.
\end{itemize}
\end{document}
